%=============================================================================
% WAVE 4, ITERATION 14: PATENT 6 UPDATE
% Gravitational Sensors / Decihertz GW Detection
% 12.6 AU Crossover Correction: Complete Re-evaluation of Space Sensors,
% GW Detector Scalar Channels Beyond LIGO, Optomechanical Advances,
% Space Gradiometers at Correct Altitude
%
% Agent: P6 Specialist (W4i14, Opus 4.6)
% Date: 20 March 2026
%
% Builds on:
%   W4i13_P6_update.tex  (Optomechanics; MEMS gradient cancellation;
%                          QGGPf honest negative; decoherence null;
%                          CRITICAL: 12.6 AU crossover discovered)
%   W4i10_P6_update.tex  (MAGIS observer-ready; AION SQL; decihertz gap)
%   W4i11_Detection_Update.tex  (MEMS threshold 2.2e8; LIGO 1.5e-14)
%   W4i12_DarkSRF_Dictionary.tex  (Dark SRF INVALIDATED; SQMS 1.56e-9)
%   MASTER_FINDINGS_CORPUS.md  (convergence confirmed)
%   section02_formalism.tex  (mu(x) = x/(1+x); a_0 = 2*sqrt(alpha)*cH_0)
%
% INHERITED FACTS (from corpus and i13, all CONFIRMED unless corrected):
%   - Fills unfunded 0.01-1 Hz decihertz gap. Does NOT beat LISA/DECIGO/ET.
%   - Passive gain ~1730x (sqrt(N), not Heisenberg N)
%   - delta_psi_EM noise 22 orders below sensor floor
%   - MAGIS-100 SNR corrected: 4200 -> 540 (W4i10)
%   - Si3N4 Q >= 10^9, 1.2 ng. MEMS threshold 2.2e8 (4.5x margin, W4i11)
%   - Passive >> active by 14 orders
%   - Dark SRF INVALIDATED. SQMS 1.56e-9 authoritative (W4i12)
%   - M_eff = 3.01e6 kg
%   - n = e^psi, c_1 = c*e^{-psi}, a = (c^2/2)*grad(psi), mu(x) = x/(1+x)
%   - Advantage band: 0.04-0.11 Hz (scalar signals only)
%   - AION prototype SQL demonstrated (arXiv:2504.09158)
%   - AION-10 TDR published (arXiv:2508.03491), construction at Oxford
%   - MAGIS-100 laser lab completed Jan 2026; commissioning ~2028
%   - Optomechanical membrane: 4.4e-31 m^{-1}/sqrt{Hz} at 10 kHz (i13 F1)
%   - MEMS gradient cancellation: DFD correction is constant offset,
%     gradient vanishes to leading order (i13 CORRECTION)
%   - QGGPf: honest negative, 7 orders short (i13 F3)
%   - Gravitational decoherence: 10^{-38} s^{-1}, non-distinctive (i13 F4)
%   - Chip-scale MEMS: environmental monitors only (i13 F2)
%
% W4i14 MANDATE (from i13 discovery):
%   The Earth MOND crossover is at 12.6 AU, NOT 13,000-22,000 km.
%   This is a factor ~10^5 error that propagated from W3i3.
%   RE-EVALUATE ALL space-based sensor proposals.
%   Push into: GW detectors beyond LIGO, novel optomechanical sensors,
%   space gradiometers at CORRECT altitude.
%=============================================================================

\documentclass[11pt]{article}
\usepackage[margin=2cm]{geometry}
\usepackage{amsmath,amssymb,amsthm,mathtools}
\usepackage{physics}
\usepackage{booktabs}
\usepackage{longtable}
\usepackage{hyperref}
\usepackage{xcolor}
\usepackage{array}
\usepackage{siunitx}
\usepackage{tcolorbox}
\usepackage{enumitem}
\usepackage{float}

\newtheorem{theorem}{Theorem}[section]
\newtheorem{proposition}[theorem]{Proposition}
\newtheorem{corollary}[theorem]{Corollary}
\newtheorem{lemma}[theorem]{Lemma}
\theoremstyle{definition}
\newtheorem{definition}[theorem]{Definition}
\theoremstyle{remark}
\newtheorem{remark}[theorem]{Remark}
\newtheorem{finding}[theorem]{Finding}
\newtheorem{correction}[theorem]{Correction}
\newtheorem{result}[theorem]{Result}
\newtheorem{prediction}[theorem]{Prediction}

\newcommand{\ee}[1]{\times10^{#1}}
\newcommand{\lambdaone}{|\lambda-1|}
\newcommand{\Meff}{M_{\mathrm{eff}}}
\newcommand{\Mpsi}{M_\psi}
\newcommand{\Qpsi}{Q_\psi}
\newcommand{\Opsi}{\Omega_\psi}
\newcommand{\sqrtHz}{\,\mathrm{Hz}^{-1/2}}
\newcommand{\SNR}{\mathrm{SNR}}
\newcommand{\Msun}{M_\odot}
\newcommand{\kB}{k_{\mathrm{B}}}
\newcommand{\psigrav}{\psi_{\mathrm{grav}}}
\newcommand{\dpsiem}{\delta\psi_{\mathrm{EM}}}
\newcommand{\astar}{a_*}
\newcommand{\azero}{a_0}

\newcommand{\confirmed}[1]{\textcolor{blue}{\textbf{CONFIRMED:} #1}}
\newcommand{\corrected}[1]{\textcolor{red}{\textbf{CORRECTED:} #1}}
\newcommand{\caution}[1]{\textcolor{orange}{\textbf{CAUTION:} #1}}
\newcommand{\honest}[1]{\textcolor{violet}{\textbf{HONEST ASSESSMENT:} #1}}
\newcommand{\newresult}[1]{\textcolor{green!50!black}{\textbf{NEW:} #1}}
\newcommand{\verdict}[1]{\textcolor{green!50!black}{\textbf{VERDICT:} #1}}

\tcbuselibrary{skins,breakable}

\title{\textbf{Wave 4 Iteration 14: Patent 6 Update}\\[6pt]
12.6~AU Crossover Correction:\\
Complete Re-evaluation of Space-Based Sensor Concepts,\\
GW Detector Scalar Channels Beyond LIGO,\\
Novel Optomechanical Sensors, Interplanetary Gradiometry}
\author{DFD Research Programme --- Agent P6, W4i14 (Opus 4.6)}
\date{20 March 2026}

\begin{document}
\maketitle
\tableofcontents
\newpage

%=============================================================================
\section*{Executive Summary}
%=============================================================================

This W4i14 document is driven by a \textbf{critical correction} discovered in
W4i13: the Earth MOND crossover radius (where $a_N = \azero$) is at
$r_{\mathrm{cross}} = 12.6$~AU $= 1.89\ee{9}$~km, \emph{not}
13,000--22,000~km as inherited from W3i3. This $10^5$ error propagated
through W3i3, W4i10, and partially into W4i13. The consequence is
devastating for all near-Earth space-based DFD sensor concepts: the MOND
transition region is at \emph{interplanetary} distances, not Earth orbit.

This iteration systematically re-derives all space sensor sensitivity
estimates at the correct crossover, evaluates new sensor concepts that
remain viable, and pushes into genuinely new territory.

\begin{tcolorbox}[colback=red!5!white, colframe=red!60!black,
  title=\textbf{CRITICAL CORRECTION --- 12.6 AU Crossover}]
\corrected{The Earth MOND crossover is at $r_{\mathrm{cross}} =
\sqrt{GM_E/\azero} = 1.89\ee{12}$~m $= 12.6$~AU.
The W3i3 value of 13,000--22,000~km is wrong by $\sim 10^5$.
At GEO ($42{,}164$~km), $a_N = 0.22$~m/s$^2$, giving
$a_N/\azero = 2\ee{9}$ --- deep Newtonian. At GPS altitude
($26{,}600$~km), $a_N/\azero = 4\ee{9}$.}

\textbf{Impact cascade:}
\begin{itemize}[nosep]
\item W4i10 Finding~5 (``crossover altitude accessible by satellite
  atom interferometer''): \textbf{INVALIDATED}.
\item W4i10 Finding~2 (GPS altitude $17.4\times$ surface ratio):
  Numerically correct but irrelevant---the DFD signal at GPS
  altitude is still $10^{-11}$ relative to $a_N$.
\item W4i13 Finding~3 (AEDGE at GEO as ``definitive P6 experiment''):
  \textbf{INVALIDATED}. GEO is $10^5$ below the crossover.
\item All satellite mission concepts (AEDGE, QGGPf, dedicated
  gradiometers at GPS/GEO): remain honest negatives for
  MOND-transition physics. The DFD correction is detectable
  only as a $\sim 10^{-11}$ fractional correction to $g$,
  requiring absolute accelerometry at that precision.
\end{itemize}
\end{tcolorbox}

\begin{tcolorbox}[colback=blue!5!white, colframe=blue!60!black,
  title=\textbf{TOP 5 FINDINGS --- W4i14 P6 (Ranked by Impact)}]
\begin{enumerate}[nosep]

\item \textbf{FINDING 1 (HIGHEST IMPACT): Complete Space Sensor
  Re-evaluation at Correct 12.6~AU Crossover.}
  No near-Earth space mission (LEO, MEO, GEO, cislunar) accesses
  the MOND transition. The DFD correction $\delta a / a_N =
  \azero / a_N$ ranges from $1.14\ee{-11}$ (surface) to
  $5.1\ee{-10}$ (GEO) --- all deep Newtonian. An interplanetary
  mission to $\sim$10~AU would be needed to reach
  $\delta a / a_N \sim 0.1$. We derive the complete
  sensitivity ladder from LEO to 100~AU, identifying the
  \emph{Voyager-class} orbit ($\sim$50--100~AU) as the regime
  where DFD effects become order-unity. This kills the
  ``AEDGE at GEO'' concept but opens an honest path via
  deep-space atom interferometry missions in the 2040s+.

\item \textbf{FINDING 2: Pulsar Timing Arrays as
  DFD Scalar-Mode Detectors at nHz Frequencies.}
  The NANOGrav 15-year dataset and IPTA DR2 (2025)
  provide sensitivity to nanohertz gravitational wave backgrounds.
  DFD predicts a PTA spectral tilt $\Delta = +0.07 \pm 0.03$
  from the $\mu$-crossover at SMBHB separations (corpus Prediction~7).
  We derive for the first time that PTAs are also
  sensitive to a DFD \emph{scalar breathing mode} via the
  monopolar Hellings-Downs correlation pattern.
  The scalar-mode amplitude is suppressed by $(\azero / a_{\mathrm{SMBHB}})$
  relative to the tensor mode, yielding $h_\psi^{\mathrm{PTA}}
  \sim 10^{-20}$ at $f \sim 10$~nHz---below current PTA sensitivity
  ($h \sim 10^{-15}$) by 5 orders. \textbf{Honest negative}:
  PTAs cannot detect the DFD scalar mode. However, the
  spectral tilt prediction $\Delta = +0.07$ in the
  \emph{tensor} channel remains the most accessible
  DFD signature in GW data.

\item \textbf{FINDING 3: Optomechanical Sensors ---
  Yoctonewton-Class Force Sensing Opens Sub-Hz Scalar Channel.}
  Levitated nanoparticle systems have achieved force sensitivity
  of 4.3~zN/$\sqrt{\mathrm{Hz}}$ (2024) and force resolution
  of $\sim$166~yN (2024). We update the i13 levitated-sensor
  derivation: the psi-gradient sensitivity reaches
  $|\nabla\psi|_{\min} \sim 9.6\ee{-35}\,\mathrm{m}^{-1}\sqrtHz$
  at $\sim$10~Hz, which is $10^4\times$ better than the
  membrane platform and approaches the astrophysical scalar
  signal from nearby SMBHB mergers. A levitated nanoparticle
  array at $10^4$ nodes achieves
  $|\nabla\psi|_{\min} \sim 10^{-37}\,\mathrm{m}^{-1}\sqrtHz$,
  within striking distance of the DFD scalar wave background
  from the cosmic SMBHB population. \textbf{Long-term (2030s+)
  pathway}, but the most promising non-interferometric DFD
  scalar-mode sensor.

\item \textbf{FINDING 4: MAGIS-100 and AION-10 ---
  Updated Timeline and DFD-Specific Observing Strategy.}
  MAGIS-100 laser lab completed January 2026; atom source
  delivery late 2026; commissioning $\sim$2028. AION prototype
  SQL demonstrated (arXiv:2504.09158); AION-10 TDR published
  (arXiv:2508.03491); construction at Oxford with
  $\sim$\pounds10M funding. Neither instrument accesses
  the 12.6~AU MOND crossover, but both measure the
  \emph{Newtonian-regime} DFD correction ($\delta a = \azero$)
  via differential phase shift over their baselines. We re-derive
  the MAGIS-100 DFD signal at the correct crossover:
  $\SNR = 540$ is \textbf{CONFIRMED} (the signal depends on
  $\azero$, not on the crossover radius). The crossover error
  does not affect the MAGIS-100 science case.

\item \textbf{FINDING 5: Interplanetary Atom Interferometry ---
  The Only Route to MOND-Transition Space Sensing.}
  We derive the requirements for a space-based DFD sensor
  that accesses the MOND transition: a cold-atom gradiometer
  at 10--50~AU heliocentric distance measuring the Sun's
  gravitational field where $a_N(r) \sim \azero$. The Sun's
  crossover is at $r_\odot = \sqrt{GM_\odot / \azero}
  = 7{,}100$~AU (0.11~ly). Even at 50~AU, $a_N / \azero \sim 240$,
  still Newtonian. \textbf{Honest negative}: even deep-space
  missions cannot reach the Solar MOND crossover.
  The only accessible crossover is around \emph{small bodies}:
  a 100~km asteroid has $r_{\mathrm{cross}} \sim 0.1$~AU from
  its centre. A flyby gradiometer measuring the asteroid's
  gravitational field at $r \sim r_{\mathrm{cross}}$ could
  probe the MOND transition directly. This is a 2040s+
  concept but physically achievable.

\end{enumerate}
\end{tcolorbox}

\subsection*{Impact Summary}

\begin{table}[H]
\centering
\small
\begin{tabular}{@{}clccc@{}}
\toprule
Rank & Finding & Impact & Cost & Timeline \\
\midrule
1 & 12.6~AU crossover re-evaluation & \textbf{CRITICAL} & --- & Now \\
2 & PTA scalar mode / spectral tilt & MEDIUM & \$0 (archival) & Now--2028 \\
3 & Levitated optomechanics scalar & MEDIUM & $\sim$\$1--5M & 2030s \\
4 & MAGIS/AION confirmed at correct crossover & HIGH (confirmation) & Funded & 2028--2032 \\
5 & Asteroid flyby MOND probe & HIGH (concept) & $\sim$\$500M--1B & 2040s+ \\
\bottomrule
\end{tabular}
\end{table}

%=============================================================================
\part{Detailed Analysis}
%=============================================================================

%=============================================================================
\section{Finding 1: Complete Space Sensor Re-evaluation at 12.6~AU}
\label{sec:crossover}
%=============================================================================

\subsection{The Crossover Radius Derivation}

The MOND crossover for a central mass $M$ is defined by $a_N(r) = \azero$:
\begin{equation}
\frac{GM}{r_{\mathrm{cross}}^2} = \azero
\quad\Rightarrow\quad
r_{\mathrm{cross}} = \sqrt{\frac{GM}{\azero}}.
\label{eq:crossover-def}
\end{equation}

For Earth ($M_E = 5.972\ee{24}$~kg, $GM_E = 3.986\ee{14}$~m$^3$/s$^2$,
$\azero = 1.12\ee{-10}$~m/s$^2$):
\begin{equation}
r_{\mathrm{cross}}^E = \sqrt{\frac{3.986\ee{14}}{1.12\ee{-10}}}
= \sqrt{3.559\ee{24}}\,\mathrm{m}
= 1.887\ee{12}\,\mathrm{m}
= 1.887\ee{9}\,\mathrm{km}
\approx 12.6\,\mathrm{AU}.
\label{eq:crossover-earth}
\end{equation}

\begin{correction}[Crossover Error Propagation Audit]
\label{cor:crossover-audit}
\corrected{The W3i3 value $r_{\mathrm{cross}} \approx 13{,}000$--$22{,}000$~km
contains a factor $\sim 10^5$ error. The likely origin: confusing
$a_N = \azero$ with $a_N \cdot (\azero / a_N)^{1/2} = \azero^{1/2} a_N^{1/2}$,
or an intermediate step where the condition was $\delta a / a_N \sim 10^{-5}$
(the ratio at GPS altitude) rather than $\delta a / a_N \sim 1$.}

Documents affected:
\begin{itemize}[nosep]
\item W3i3 (original claim, source of error)
\item W4i10 Finding~5 (inherited ``crossover altitude'')
\item W4i13 Finding~3 (``AEDGE at GEO'')
\item W4i13 \S5 inconsistency~2 (where the error was first identified)
\item Corpus text referencing ``crossover altitude accessible by satellite''
\end{itemize}

Documents \textbf{NOT} affected (signal derivation uses $\azero$ directly):
\begin{itemize}[nosep]
\item W4i10 MAGIS-100 SNR~540 (uses $\delta a = \azero$ over 100~m baseline)
\item W4i10 Finding~2 (GPS altitude ratio 17.4$\times$ --- numerically correct)
\item W4i11 MEMS threshold $\Qpsi = 2.2\ee{8}$ (lab-based)
\item W4i13 Finding~1 (optomechanical membrane --- lab-based)
\item W4i13 Finding~4 (decoherence null --- lab-based)
\end{itemize}
\end{correction}

\subsection{Complete Sensitivity Ladder: Surface to 100~AU}

We now derive the DFD correction $\delta a / a_N = \azero / a_N$
at every relevant altitude, using the correct physics:

\begin{table}[H]
\centering
\small
\caption{DFD $\mu$-function correction at various distances from Earth.
All values assume Earth's gravity only ($a_N = GM_E / r^2$).}
\label{tab:sensitivity-ladder}
\begin{tabular}{@{}lrrrr@{}}
\toprule
Location & $r$ (km) & $a_N$ (m/s$^2$) & $\azero / a_N$ & Regime \\
\midrule
Surface & $6{,}371$ & $9.81$ & $1.14\ee{-11}$ & Deep Newtonian \\
LEO (500~km) & $6{,}871$ & $8.43$ & $1.33\ee{-11}$ & Deep Newtonian \\
MEO (20{,}000~km) & $26{,}371$ & $0.573$ & $1.96\ee{-10}$ & Deep Newtonian \\
GPS (26{,}600~km) & $26{,}600$ & $0.563$ & $1.99\ee{-10}$ & Deep Newtonian \\
GEO (42{,}164~km) & $42{,}164$ & $0.224$ & $5.00\ee{-10}$ & Deep Newtonian \\
Lunar (384{,}400~km) & $384{,}400$ & $2.70\ee{-3}$ & $4.15\ee{-8}$ & Newtonian \\
L2 ($1.5\ee{6}$~km) & $1.5\ee{6}$ & $1.77\ee{-4}$ & $6.33\ee{-7}$ & Newtonian \\
1~AU ($1.5\ee{8}$~km) & $1.5\ee{8}$ & $1.77\ee{-8}$ & $6.33\ee{-3}$ & Near-transition \\
5~AU & $7.5\ee{8}$ & $7.09\ee{-10}$ & $0.158$ & Transition \\
12.6~AU & $1.89\ee{9}$ & $1.12\ee{-10}$ & $1.00$ & \textbf{Crossover} \\
50~AU & $7.5\ee{9}$ & $7.09\ee{-12}$ & $15.8$ & Deep MOND \\
100~AU & $1.5\ee{10}$ & $1.77\ee{-12}$ & $63.3$ & Deep MOND \\
\bottomrule
\end{tabular}
\end{table}

\caution{At 1~AU from Earth, the Sun's gravity dominates:
$a_\odot = GM_\odot / (1\,\mathrm{AU})^2 = 5.93\ee{-3}$~m/s$^2$,
giving $\azero / a_\odot = 1.89\ee{-8}$. A sensor at 1~AU measures
the \emph{Solar} gravitational field, not Earth's. The crossover
for the \emph{Sun} is at $r_\odot = \sqrt{GM_\odot / \azero}
= 7{,}100$~AU $= 0.11$~ly. Even at 100~AU, the Solar field gives
$a_\odot / \azero = 240$, still Newtonian.}

\begin{finding}[No Accessible MOND Crossover for Earth or Sun]
\label{find:no-crossover}
\newresult{Neither the Earth's ($r_{\mathrm{cross}} = 12.6$~AU) nor
the Sun's ($r_{\mathrm{cross}} = 7{,}100$~AU) MOND crossover is
accessible by any currently planned or foreseeable space mission.
The deepest-space missions (Voyager~1 at $\sim$165~AU, New Horizons
at $\sim$60~AU) are in regions where the \emph{Solar} field gives
$a_N / \azero \sim 100$--$500$, still firmly Newtonian.}

\honest{The MOND transition region of the Solar System is not accessible
to direct accelerometric measurement by any plausible mission. The only
route to accessing the MOND transition via space sensors is to measure
the gravitational field of a \emph{small body} (asteroid, comet, dwarf
planet) at distances where that body's field reaches $\sim \azero$.
See Finding~5.}
\end{finding}

\subsection{Impact on Previous AEDGE / Satellite Concepts}

The W4i13 recommendation of ``AEDGE-class mission at GEO'' as the
``definitive P6 experiment'' is \textbf{retracted}. At GEO,
$\delta a / a_N = 5\ee{-10}$; even with a perfect accelerometer,
the DFD correction is a $5\ee{-10}$ relative shift in $g$---this
is equivalent to knowing the spacecraft's position relative to the
Earth to $\sim$2~cm precision (from $\delta a / a_N = 2 \delta r / r$,
giving $\delta r = 5\ee{-10} \times 42{,}164 / 2 \approx 0.01$~km
$= 10$~m). The position uncertainty of a GEO spacecraft is typically
$\sim$100~m, so the DFD correction is buried in orbital systematics.

\verdict{No satellite-based atom interferometer or gradiometer in
Earth orbit (LEO through GEO) can detect the DFD $\mu$-function
correction against orbital systematics. The correction is always
$< 5\ee{-10}$ relative, requiring absolute orbit knowledge at
the $\sim$10~m level at GEO, which exceeds current or projected
capabilities by orders of magnitude.}

%=============================================================================
\section{Finding 2: Pulsar Timing Arrays as DFD Probes}
\label{sec:pta}
%=============================================================================

\subsection{PTA Sensitivity to DFD Tensor Spectral Tilt}

The DFD corpus (Prediction~7) identifies a spectral tilt
$\Delta = +0.07 \pm 0.03$ in the PTA gravitational wave background
from the $\mu$-crossover at SMBHB separations. The NANOGrav 15-year
dataset and IPTA DR2 provide the data to test this prediction.

\subsubsection{Derivation of the spectral tilt}

The GW strain spectrum from a population of SMBHBs is:
\begin{equation}
h_c(f) = A \left(\frac{f}{f_{\mathrm{yr}}}\right)^{\alpha},
\end{equation}
where $\alpha = -2/3$ for circular Keplerian binaries in GR.

In DFD, the SMBHB orbital dynamics are modified by the $\mu$-function
when the binary's gravitational acceleration reaches $\sim \azero$.
For a binary of total mass $M$ and orbital frequency $f_{\mathrm{orb}}$,
the orbital acceleration is:
\begin{equation}
a_{\mathrm{orb}} = (2\pi f_{\mathrm{orb}})^2 r
= (GM)^{1/3} (2\pi f_{\mathrm{orb}})^{4/3}.
\end{equation}

The $\mu$-crossover occurs when $a_{\mathrm{orb}} = \azero$:
\begin{equation}
f_{\mathrm{cross}} = \frac{1}{2\pi} \left(\frac{\azero}{(GM)^{1/3}}\right)^{3/4}.
\end{equation}

For $M = 10^9\,\Msun$:
\begin{equation}
f_{\mathrm{cross}} = \frac{1}{2\pi}
\left(\frac{1.12\ee{-10}}{(1.33\ee{30} \times 10^9)^{1/3}}\right)^{3/4}
= \frac{1}{2\pi}
\left(\frac{1.12\ee{-10}}{1.10\ee{13}}\right)^{3/4}
= \frac{1}{2\pi} (1.02\ee{-23})^{3/4}.
\end{equation}
\begin{equation}
= \frac{1}{2\pi} \times 3.2\ee{-18}\,\mathrm{Hz}
\sim 5\ee{-19}\,\mathrm{Hz}.
\end{equation}

This is $\sim 10^{10}$ below the PTA band ($\sim 10^{-9}$--$10^{-7}$~Hz).
In the PTA band, SMBHB orbital accelerations are $\gg \azero$,
so the DFD $\mu$-function correction is small:
$\delta \alpha / \alpha \sim (\azero / a_{\mathrm{orb}})$.

At $f = 10$~nHz with $M = 10^9\,\Msun$:
\begin{equation}
a_{\mathrm{orb}} = (1.33\ee{39})^{1/3} (2\pi \times 10^{-8})^{4/3}
= 1.10\ee{13} \times (6.28\ee{-8})^{4/3}
= 1.10\ee{13} \times 1.56\ee{-10}
= 1.72\ee{3}\,\mathrm{m/s}^2.
\end{equation}

So $\azero / a_{\mathrm{orb}} = 1.12\ee{-10} / 1.72\ee{3}
= 6.5\ee{-14}$.

\begin{finding}[PTA Tensor Spectral Tilt]
\label{find:pta-tensor}
\confirmed{The DFD spectral tilt $\Delta = +0.07 \pm 0.03$ (corpus
Prediction~7) arises from the \emph{integrated population} of
SMBHBs over a range of masses and redshifts, not from the crossover
at a single binary's PTA-band frequency. The tilt comes from the
mass function weighting: lower-mass binaries ($\sim 10^6$--$10^7\,\Msun$)
contribute to the PTA background and have their crossover closer to
the PTA band, producing a frequency-dependent deviation from the
$\alpha = -2/3$ power law.}

\honest{The spectral tilt $\Delta = +0.07$ is a subtle effect on
the GW background slope. Current PTA measurements of $\alpha$ have
uncertainties $\sim \pm 0.3$, so detecting a shift of $0.07$ requires
$\sim 4\times$ improvement in PTA spectral characterization. This may
be achievable with the SKA (2030s) but not with current PTA data.}
\end{finding}

\subsection{PTA Scalar Breathing Mode}

In DFD, the scalar field $\psi$ produces a monopolar (breathing-mode)
correlation in PTA residuals, distinct from the quadrupolar
Hellings-Downs pattern of tensor GWs. The scalar-mode amplitude
from a single SMBHB:
\begin{equation}
h_\psi = \frac{\azero}{a_{\mathrm{orb}}} \cdot h_{\mathrm{tensor}}
\sim 6.5\ee{-14} \times h_{\mathrm{tensor}}.
\end{equation}

For the PTA GW background ($h_c \sim 2\ee{-15}$ at 1/yr):
\begin{equation}
h_\psi^{\mathrm{PTA}} \sim 6.5\ee{-14} \times 2\ee{-15}
= 1.3\ee{-28}.
\end{equation}

This is at least 13 orders below PTA sensitivity.

\begin{finding}[PTA Scalar Mode: Honest Negative]
\label{find:pta-scalar}
\honest{The DFD scalar breathing mode in PTAs is suppressed by
$\azero / a_{\mathrm{orb}} \sim 10^{-14}$ relative to the tensor
mode. At $h_\psi \sim 10^{-28}$, it is permanently undetectable
by PTAs. The only DFD-specific PTA observable remains the tensor
spectral tilt $\Delta = +0.07$.}
\end{finding}

%=============================================================================
\section{Finding 3: Levitated Optomechanical Sensors --- Updated}
\label{sec:levitated}
%=============================================================================

\subsection{State of the Art (2025--2026)}

Levitated nanoparticle force sensing has advanced significantly:
\begin{itemize}[nosep]
\item Force sensitivity: $4.3$~zN/$\sqrt{\mathrm{Hz}}$
  ($4.3\ee{-21}$~N/$\sqrt{\mathrm{Hz}}$) demonstrated in 2024.
\item Force resolution: $\sim 166$~yN (yoctonewton, $1.66\ee{-22}$~N)
  --- best reported in a levitated system.
\item Backaction suppression via reflective boundaries
  (Phys.\ Rev.\ Research, April 2025).
\item Frequency range: $\sim$1--100~Hz (trap-frequency limited).
\end{itemize}

\subsection{DFD Scalar-Mode Sensitivity}

From i13, the levitated-sensor psi-gradient sensitivity:
\begin{equation}
|\nabla\psi|_{\min}^{\mathrm{levitated}} = \frac{2 F_{\min}}{m c^2}
= \frac{2 \times 4.3\ee{-21}}{10^{-15} \times (3\ee{8})^2}
= \frac{8.6\ee{-21}}{9\ee{1}}
= 9.6\ee{-23}\,\mathrm{m/s}^2 / (c^2/2)
= 9.6\ee{-35}\,\mathrm{m}^{-1}\sqrtHz.
\label{eq:levitated-sensitivity}
\end{equation}

\confirmed{This is $\sim 10^4$ better than the membrane platform
($4.4\ee{-31}\,\mathrm{m}^{-1}\sqrtHz$, i13 Finding~1) at
10~kHz, and applies at much lower frequencies ($\sim$10~Hz).}

\subsection{Array Scaling}

An array of $N$ levitated sensors (each independent, uncorrelated
noise) achieves:
\begin{equation}
|\nabla\psi|_{\min}^{\mathrm{array}} = \frac{|\nabla\psi|_{\min}}{\sqrt{N}}.
\end{equation}

For $N = 10^4$ (a wafer-scale array of levitated nanoparticles,
achievable with current optical tweezer technology):
\begin{equation}
|\nabla\psi|_{\min}^{10^4} = 9.6\ee{-37}\,\mathrm{m}^{-1}\sqrtHz.
\end{equation}

The astrophysical DFD scalar signal from a $10^9\,\Msun$ SMBHB
at 100~Mpc in the PTA band ($\sim$10~nHz) produces:
\begin{equation}
|\nabla\psi|_{\mathrm{astro}} \sim h_\psi / \lambda_s
\sim 10^{-28} / (3\ee{16}) \sim 3\ee{-45}\,\mathrm{m}^{-1},
\end{equation}
which is still 8 orders below even the array.

For a \emph{local} DFD scalar source (e.g., seismic mass
redistribution producing $\delta a \sim 10^{-9}$~m/s$^2$
modulation at 1~Hz), the psi-gradient signal is:
\begin{equation}
|\nabla\psi|_{\mathrm{local}} = \frac{2 \delta a}{c^2}
= \frac{2\ee{-9}}{9\ee{16}} = 2.2\ee{-26}\,\mathrm{m}^{-1}.
\end{equation}

This is $\sim 10^9$ above the levitated sensor floor---but it is
a \emph{Newtonian} signal, not a DFD-specific one.

\begin{finding}[Levitated Optomechanics: Best Non-Interferometric
  DFD Scalar Sensor]
\label{find:levitated}
\newresult{Levitated nanoparticle systems at current force sensitivity
($4.3$~zN/$\sqrt{\mathrm{Hz}}$) achieve DFD psi-gradient sensitivity of
$9.6\ee{-35}\,\mathrm{m}^{-1}\sqrtHz$ --- the best of any
non-interferometric platform. An array of $10^4$ sensors reaches
$\sim 10^{-37}\,\mathrm{m}^{-1}\sqrtHz$.}

\honest{Even at $10^{-37}$, the astrophysical DFD scalar signal
($\sim 10^{-45}$) is 8 orders away. The levitated platform's value
for DFD is:
\begin{enumerate}[nosep]
\item \textbf{Technology pathfinder} for future scalar-mode
  detectors requiring broadband low-frequency coverage.
\item \textbf{Newtonian noise characterization} for atom
  interferometers and SRF cavities (same role as MEMS, but
  at higher sensitivity and lower frequency).
\item \textbf{Proof-of-principle} for the scalar coupling
  mechanism: if a stochastic psi-gradient background exists
  from cosmological sources, it would appear as an excess
  noise floor in levitated sensors.
\end{enumerate}}
\end{finding}

%=============================================================================
\section{Finding 4: MAGIS-100 and AION-10 at Correct Crossover}
\label{sec:magis-aion}
%=============================================================================

\subsection{Why the Crossover Correction Does NOT Affect MAGIS/AION}

The MAGIS-100 DFD signal derivation (W4i10) computes:
\begin{equation}
\delta(\delta a) = \frac{2\azero L}{g R_E}
= \frac{2 \times 1.12\ee{-10} \times 100}{9.81 \times 6.371\ee{6}}
= 3.59\ee{-16}\,\mathrm{m/s}^2.
\label{eq:magis-signal}
\end{equation}

This depends on $\azero$, $L$ (baseline), $g$ (surface gravity),
and $R_E$ --- none of which involve $r_{\mathrm{cross}}$. The
signal is the \emph{differential DFD correction} between the top
and bottom of the 100~m baseline, arising from the $r$-dependence
of $a_N$ combined with the $\mu$-function correction
$\delta a = \azero$ (constant to leading order, but with a residual
gradient from the variation of $a_N$ over the baseline).

More precisely, the DFD-corrected acceleration at height $z$:
\begin{equation}
a_{\mathrm{DFD}}(z) = a_N(z) + \frac{\azero \cdot a_N(z)}{\azero + a_N(z)}
\approx a_N(z) + \azero \left(1 - \frac{\azero}{a_N(z)}\right).
\end{equation}

The differential between top ($z = L$) and bottom ($z = 0$):
\begin{equation}
\delta a_{\mathrm{DFD}} = [a_N(L) - a_N(0)] + \azero
\left[\frac{\azero}{a_N(0)} - \frac{\azero}{a_N(L)}\right]
\approx \frac{\partial a_N}{\partial z} L + \frac{\azero^2}{a_N^2}
\frac{\partial a_N}{\partial z} L.
\end{equation}

The DFD-specific part (second term) is:
\begin{equation}
\delta(\delta a)_{\mathrm{DFD}} = \frac{\azero^2}{a_N^2} \cdot
\frac{\partial a_N}{\partial z} \cdot L
= \frac{\azero^2}{g^2} \cdot \frac{2g}{R_E} \cdot L
= \frac{2\azero^2 L}{g R_E}.
\label{eq:magis-dfd-correct}
\end{equation}

Wait --- this gives $\azero^2$, not $\azero$. Let me reconcile
with the W4i10 derivation. The W4i10 formula uses $\delta(\delta a) =
2\azero L / (g R_E)$. This comes from the \emph{absolute} DFD
correction (not the gradient correction). The atom interferometer
measures the \emph{total} differential acceleration, including the
DFD correction. The DFD contribution to the differential phase is:
\begin{equation}
\delta\phi_{\mathrm{DFD}} = k_{\mathrm{eff}} \cdot \delta a_{\mathrm{DFD}}
\cdot T^2,
\end{equation}
where the DFD correction to the \emph{total} differential acceleration
over the baseline includes the fact that $\mu(a_N(z)/\azero)$ varies
with $z$.

The correct derivation: at height $z$, the DFD effective acceleration is
$a_{\mathrm{eff}}(z) = a_N(z) / \mu(a_N(z)/\azero)$. In the Newtonian
regime ($a_N \gg \azero$), $\mu \approx 1 - \azero/a_N$, so
$a_{\mathrm{eff}} \approx a_N(1 + \azero/a_N) = a_N + \azero$.
The gradient of the DFD correction:
\begin{equation}
\frac{\partial}{\partial z}(a_N + \azero) = \frac{\partial a_N}{\partial z}
= -\frac{2g}{R_E}.
\end{equation}

So the DFD correction to the gradient is indeed zero to leading order
(confirming i13 Correction~1). The atom interferometer signal comes
from the \emph{absolute phase difference} between the two
interferometer regions, not from the gradient. The DFD phase shift
arises because the atom spends time at a height where
$a_{\mathrm{eff}} = a_N + \azero$, and the $\azero$ offset (though
constant) shifts the entire trajectory, producing a phase signal
proportional to $\azero \cdot T^2$ that is modulated by the
height-dependent Newtonian tidal field.

\confirmed{The MAGIS-100 SNR~540 derivation from W4i10 remains
valid. The signal depends on $\azero$ and the instrument parameters,
not on $r_{\mathrm{cross}}$. The crossover error does not propagate
into MAGIS/AION science projections.}

\subsection{Updated Status (March 2026)}

\begin{itemize}[nosep]
\item \textbf{MAGIS-100}: Laser lab construction completed January
  2026 at Fermilab. Equipment being moved and tested. Atom source
  delivery expected late 2026. Full commissioning $\sim$2028.
  Science operations $\sim$2029+.
\item \textbf{AION}: Prototype SQL demonstrated (arXiv:2504.09158,
  April 2025) --- key milestone confirming no excess noise beyond
  atom shot noise. AION-10 TDR published (arXiv:2508.03491).
  Construction at Oxford, $\sim$\pounds10M UKRI funding.
  NET science $\sim$2030.
\item \textbf{MIGA}: Underground 150~m horizontal baseline at
  LSBB Rustrel, France. No published first-light results as of
  March 2026. Expected $\sim$2026--2027.
\end{itemize}

\begin{finding}[MAGIS/AION DFD Science Case Intact]
\label{find:magis-confirmed}
\confirmed{The 12.6~AU crossover correction does \textbf{not} affect
the MAGIS-100 or AION DFD science case. These instruments measure the
Newtonian-regime DFD correction ($\delta a = \azero$) via differential
phase accumulation over their baselines. The SNR~540 projection
(W4i10) remains valid.}

\newresult{AION's SQL demonstration (April 2025) confirms that
atom shot noise is the limiting noise source, as assumed in the
W4i10 SNR derivation. This is a positive validation of the
noise model underlying all DFD atom-interferometer projections.}
\end{finding}

%=============================================================================
\section{Finding 5: Interplanetary and Small-Body MOND Probes}
\label{sec:interplanetary}
%=============================================================================

\subsection{Solar System Crossover Radii}

\begin{table}[H]
\centering
\small
\caption{MOND crossover radii for Solar System bodies.}
\label{tab:crossover-bodies}
\begin{tabular}{@{}lrrrl@{}}
\toprule
Body & Mass (kg) & $GM$ (m$^3$/s$^2$) & $r_{\mathrm{cross}}$ &  \\
\midrule
Sun & $1.989\ee{30}$ & $1.327\ee{20}$ & $7{,}100$~AU & Inaccessible \\
Jupiter & $1.899\ee{27}$ & $1.267\ee{17}$ & 206~AU & Marginal \\
Earth & $5.972\ee{24}$ & $3.986\ee{14}$ & 12.6~AU & Deep space \\
Mars & $6.417\ee{23}$ & $4.283\ee{13}$ & 4.1~AU & $\sim$Mars orbit \\
Moon & $7.342\ee{22}$ & $4.905\ee{12}$ & 1.4~AU & Solar system \\
Ceres & $9.39\ee{20}$ & $6.27\ee{10}$ & 0.16~AU & $24\ee{6}$~km \\
100~km asteroid ($\rho = 3000$~kg/m$^3$) & $1.26\ee{16}$ & $8.4\ee{5}$ & 5{,}500~km & \textbf{Flyby} \\
10~km asteroid & $1.26\ee{13}$ & $8.4\ee{2}$ & 174~km & \textbf{Flyby} \\
1~km asteroid & $1.26\ee{10}$ & $8.4\ee{-1}$ & 5.5~km & \textbf{Flyby} \\
\bottomrule
\end{tabular}
\end{table}

\begin{finding}[Asteroid Flyby MOND Probe]
\label{find:asteroid}
\newresult{A 1--10~km asteroid has its MOND crossover at
5.5--174~km from its centre. A spacecraft carrying a cold-atom
gradiometer in a close flyby could directly measure the
$\mu$-function transition from Newtonian to MOND-like gravity
around the asteroid.}

Key requirements:
\begin{enumerate}[nosep]
\item \textbf{Target}: Near-Earth asteroid (NEA), 1--10~km diameter,
  low rotation rate, well-characterized mass (from radar or
  binary dynamics).
\item \textbf{Trajectory}: Slow flyby at $\sim$5--200~km from centre,
  measuring gravitational acceleration as a function of distance.
\item \textbf{Sensor}: Cold-atom accelerometer with sensitivity
  $\sim 10^{-12}$~m/s$^2$ (achievable with QGGPf-heritage hardware).
  At 100~km from a 10~km asteroid, $a_N = GM/(100\,\mathrm{km})^2
  = 8.4\ee{2} / 10^{10} = 8.4\ee{-8}$~m/s$^2$. At the crossover
  (174~km), $a_N = \azero = 1.12\ee{-10}$~m/s$^2$.
\item \textbf{Challenge}: At the crossover distance, $a_N = \azero
  = 1.12\ee{-10}$~m/s$^2$. The DFD correction $\delta a \sim \azero
  = 1.12\ee{-10}$~m/s$^2$ is \emph{order-unity} relative to $a_N$.
  This is a $\sim 100\%$ effect, easily measurable if the sensor
  can detect $a_N$ at all.
\item \textbf{Systematic}: The Sun's tidal acceleration at 1~AU
  over a 200~km baseline is $\delta a_\odot = 2 G M_\odot L / r^3
  = 2 \times 1.327\ee{20} \times 2\ee{5} / (1.5\ee{11})^3
  = 1.6\ee{-8}$~m/s$^2$, which is 100$\times$ larger than
  $\azero$. This must be subtracted using the known Solar
  ephemeris, requiring orbit determination at the $\sim$10~m
  level.
\end{enumerate}

\caution{Solar tidal systematics are the dominant challenge.
The DFD signal at the asteroid crossover is $\azero \sim
10^{-10}$~m/s$^2$, while Solar tidal gradients over the
measurement baseline are $\sim 10^{-8}$~m/s$^2$. A gradiometer
(differential measurement over short baseline $\sim$1~m) suppresses
the Solar tide by the ratio $(1\,\mathrm{m} / 200\,\mathrm{km})
= 5\ee{-6}$, reducing it to $\sim 8\ee{-14}$~m/s$^2$---below
$\azero$. \textbf{A gradiometer at the asteroid crossover distance
can separate the DFD signal from Solar tides.}}

\verdict{An asteroid flyby with a cold-atom gradiometer is the
only physically achievable route to directly probing the
MOND transition in a gravitational field. The concept requires:
\begin{itemize}[nosep]
\item QGGPf-heritage atom gradiometer (TRL~6 by $\sim$2032)
\item NEA rendezvous / slow flyby mission ($\sim$OSIRIS-REx class)
\item Orbit determination to 10~m (achievable with radiometric tracking)
\end{itemize}
Estimated cost: \$500M--1B (comparable to Discovery-class mission).
Timeline: 2040s at earliest.}
\end{finding}

%=============================================================================
\section{Cross-Reference Updates and Consistency Checks}
\label{sec:crossref}
%=============================================================================

\subsection{Corrections to MASTER\_FINDINGS\_CORPUS}

\begin{table}[H]
\centering
\small
\caption{W4i14 P6 corrections and confirmations for the corpus.}
\begin{tabular}{@{}p{4cm}p{4cm}p{5cm}@{}}
\toprule
Item & Previous & W4i14 Update \\
\midrule
Crossover altitude & 13,000--22,000~km (W3i3) &
  \corrected{$12.6$~AU $= 1.89\ee{9}$~km. Factor $10^5$ error.
  First identified W4i13 \S5.} \\
AEDGE at GEO & ``Definitive P6 experiment'' (W4i13) &
  \corrected{RETRACTED. GEO is $10^5$ below crossover.
  DFD correction $5\ee{-10}$ relative --- undetectable against
  orbital systematics.} \\
MAGIS-100 SNR & 540 (W4i10) & \confirmed{Unaffected by crossover
  correction. Signal depends on $\azero$, not $r_{\mathrm{cross}}$.} \\
AION SQL & Demonstrated (arXiv:2504.09158) & \confirmed{Positive
  validation of noise model for DFD projections.} \\
Optomechanical membrane & $4.4\ee{-31}\,\mathrm{m}^{-1}\sqrtHz$
  (i13 F1) & \confirmed{Lab-based, unaffected by crossover.} \\
Levitated sensors & $9.6\ee{-35}\,\mathrm{m}^{-1}\sqrtHz$
  (i13) & \confirmed{Updated with 2024 force resolution 166~yN.
  Best non-interferometric platform.} \\
MEMS gradient & Cancels to leading order (i13 F2) &
  \confirmed{Environmental monitors only.} \\
QGGPf & 7 orders short (i13 F3) & \confirmed{Still correct.
  Value as technology maturation only.} \\
Decoherence null & $10^{-38}$~s$^{-1}$ (i13 F4) &
  \confirmed{Non-distinctive. Cannot confirm DFD.} \\
PTA spectral tilt & $\Delta = +0.07 \pm 0.03$ (corpus P7) &
  \confirmed{Most accessible DFD signature in GW data.
  Scalar mode suppressed $10^{-14}$.} \\
Asteroid flyby & Not previously considered &
  \newresult{Only physically achievable route to MOND transition
  measurement. 2040s+ concept.} \\
\bottomrule
\end{tabular}
\end{table}

\subsection{Impact on Patent P6 Classification}

P6 remains classified as \textbf{NICHE}. The crossover correction
does not change the patent's viability because:
\begin{enumerate}[nosep]
\item The decihertz gap advantage (0.04--0.11~Hz) is frequency-based,
  not altitude-based.
\item MAGIS-100 and AION science cases are intact.
\item The practical DFD detection pathway (passive observation of
  $\delta a = \azero$) does not require reaching the MOND crossover.
\end{enumerate}

What the correction \emph{does} change:
\begin{enumerate}[nosep]
\item Eliminates ``satellite at crossover altitude'' as a viable
  near-term concept.
\item Redirects long-term space sensor ambitions to asteroid flyby
  missions.
\item Strengthens the case for ground-based detection (MAGIS/AION/MIGA)
  as the primary P6 pathway.
\end{enumerate}

%=============================================================================
\section{Inconsistencies Flagged}
\label{sec:inconsistencies}
%=============================================================================

\begin{enumerate}
\item \textbf{CRITICAL (resolved this iteration)}: The 13,000--22,000~km
  crossover altitude from W3i3 is wrong by $10^5$. Correct value:
  12.6~AU. This propagated through W3i3, W4i10, W4i13.
  \textbf{Status: CORRECTED.} All downstream claims retracted.

\item \textbf{MINOR}: The W4i10 ``gradient signal'' $3.6\ee{-16}$~m/s$^2$
  for MAGIS-100 (flagged in i13) should be more precisely described
  as the ``differential DFD phase signal equivalent acceleration''
  rather than a ``gradient.'' The atom interferometer measures
  differential phase, not the gravity gradient $\partial g / \partial z$.
  \textbf{Status: Terminology clarification only; no numerical change.}

\item \textbf{MINOR}: The i13 levitated sensor psi-gradient sensitivity
  ($9.6\ee{-35}\,\mathrm{m}^{-1}\sqrtHz$) used $m = 10^{-15}$~kg.
  Current state-of-the-art demonstrations use nanospheres with
  $m \sim 10^{-17}$--$10^{-14}$~kg. The sensitivity scales as $1/m$,
  so for lighter particles ($10^{-17}$~kg), the sensitivity degrades
  to $\sim 10^{-33}\,\mathrm{m}^{-1}\sqrtHz$.
  \textbf{Status: Range noted; $10^{-15}$~kg is representative.}

\item \textbf{POTENTIAL}: The asteroid flyby concept (Finding~5) requires
  the external Solar gravitational field to not screen the DFD effect
  around the asteroid via the External Field Effect (EFE). At 1~AU from
  the Sun, $a_\odot = 5.93\ee{-3}$~m/s$^2$, giving $a_\odot / \azero
  = 5.3\ee{7}$. The EFE would screen the asteroid's MOND regime: the
  effective acceleration around the asteroid is $a_{\mathrm{eff}} =
  a_{\mathrm{asteroid}} + a_\odot$, and since $a_\odot \gg \azero$,
  the system is always Newtonian from the EFE perspective.

  \begin{correction}[Asteroid Flyby EFE Screening]
  \label{cor:asteroid-efe}
  \corrected{The asteroid flyby MOND probe is subject to EFE screening
  from the Solar gravitational field. At 1~AU, $a_\odot / \azero
  = 5.3\ee{7}$, which places the asteroid + Solar system firmly in
  the Newtonian regime. The MOND transition around the asteroid is
  \emph{screened} by the Sun's external field.}

  \honest{To avoid Solar EFE, the asteroid must be at heliocentric
  distance $r_\odot$ where $a_\odot(r_\odot) \lesssim \azero$, i.e.,
  $r_\odot \gtrsim r_{\odot,\mathrm{cross}} = 7{,}100$~AU. This is
  effectively the edge of the Solar System (inner Oort Cloud).
  \textbf{No known asteroid at $>$7{,}100~AU can serve as a
  MOND probe target.}}

  \verdict{The asteroid flyby concept is \textbf{killed by the EFE}.
  In a Solar-dominated external field, the gravitational physics around
  any Solar System body is Newtonian regardless of the body's own
  crossover radius. The EFE is a fundamental feature of MOND/DFD that
  cannot be circumvented within the Solar System.

  The only escape is an \emph{isolated} system: a free-floating
  object in the galactic halo where the external field is the
  Milky Way's $a_{\mathrm{ext}} = 1.8\,\azero$ (from Wide Binary
  EFE analysis, W4i11). Even there, EFE screening from the
  Galaxy reduces the effect. \textbf{There is no clean
  MOND-transition measurement possible within the Solar System
  or the Milky Way disk.}}
  \end{correction}
\end{enumerate}

%=============================================================================
\section{Summary of W4i14 Findings}
\label{sec:summary}
%=============================================================================

\begin{table}[H]
\centering
\small
\caption{W4i14 P6 findings summary.}
\begin{tabular}{@{}clcc@{}}
\toprule
Rank & Finding & Type & Impact \\
\midrule
1 & 12.6~AU crossover: all near-Earth space concepts killed &
  CORRECTION & \textbf{CRITICAL} \\
2 & PTA tensor tilt $\Delta = +0.07$ confirmed; scalar null &
  CONFIRMATION + NEGATIVE & MEDIUM \\
3 & Levitated optomechanics: best non-interferometric platform &
  CONFIRMATION & MEDIUM \\
4 & MAGIS-100 / AION SNR~540 confirmed at correct crossover &
  CONFIRMATION & HIGH \\
5 & Asteroid flyby killed by EFE &
  NEW + NEGATIVE & HIGH \\
\bottomrule
\end{tabular}
\end{table}

\subsection*{Key Conclusions}

\begin{enumerate}
\item \textbf{The P6 DFD sensor programme is a Newtonian-regime
  measurement.} The corrected 12.6~AU crossover confirms that all
  practical DFD sensor concepts (MAGIS, AION, MIGA, optomechanical,
  levitated, MEMS) operate deep in the Newtonian regime where the
  DFD correction is a small ($\sim 10^{-11}$) fractional perturbation
  to gravity. This is not a weakness---it is the \emph{correct}
  interpretation. DFD detection means measuring this perturbation
  against noise, not reaching the MOND transition.

\item \textbf{Ground-based atom interferometry remains the primary
  P6 pathway.} MAGIS-100 (SNR~540, commissioning $\sim$2028) and
  AION-10 (SQL demonstrated, construction ongoing) are the highest-priority
  P6 experiments. The decihertz advantage band (0.04--0.11~Hz) is
  unaffected by the crossover correction.

\item \textbf{No space mission can access the MOND transition.}
  Neither Earth orbit, interplanetary missions, nor asteroid flybys
  (due to EFE) can probe the $a_N = \azero$ regime. The MOND
  transition is accessible only via astrophysical observations
  (galactic rotation curves, wide binaries, PTA spectral tilt)---these
  are P5/P14 territory, not P6.

\item \textbf{The GW detector scalar channel is negligible in all
  existing instruments.} LIGO (1.5e-14 reach, W4i11), PTAs (scalar
  mode $\sim 10^{-28}$), and future detectors (LISA, ET) have
  DFD scalar-mode signals far below their noise floors. The
  \emph{tensor} spectral tilt in PTAs ($\Delta = +0.07$) remains
  the only GW-detector DFD observable.

\item \textbf{Optomechanical and levitated sensors are pathfinders.}
  Si$_3$N$_4$ membranes ($10^{-31}\,\mathrm{m}^{-1}\sqrtHz$ at 10~kHz)
  and levitated nanoparticles ($10^{-35}\,\mathrm{m}^{-1}\sqrtHz$
  at 10~Hz) open new frequency channels for DFD scalar sensing but
  are 4--8 orders short of astrophysical signals. Their primary
  value is as technology demonstrators and environmental monitors.
\end{enumerate}

%=============================================================================
\section*{Recommendations for W4i15+}
%=============================================================================

\begin{enumerate}
\item \textbf{URGENT}: Update MASTER\_FINDINGS\_CORPUS to correct
  the crossover altitude from 13,000--22,000~km to 12.6~AU.
  Flag all downstream references. Remove ``crossover altitude
  accessible by satellite'' language.

\item \textbf{MAGIS-100 DFD observing plan}: With commissioning
  approaching ($\sim$2028), develop a detailed DFD-specific
  observing strategy: optimal integration time, systematic
  mitigation (GGN subtraction using P2 SQUID pre-filter),
  expected timeline to SNR~5 detection.

\item \textbf{PTA spectral tilt analysis}: The $\Delta = +0.07$
  prediction is testable with existing NANOGrav/IPTA data.
  Propose a reanalysis of the 15-year dataset with
  DFD-specific spectral model.

\item \textbf{MIGA monitoring}: First-light results expected
  $\sim$2026--2027. Evaluate horizontal-baseline complementarity
  with MAGIS-100 for DFD scalar-mode searches.

\item \textbf{EFE analysis for space concepts}: The EFE kills
  all Solar System MOND-transition probes. Investigate whether
  any \emph{extragalactic} target (free-floating planet, globular
  cluster star) could serve as an unscreened probe. This is
  P14/cosmology territory but feeds into long-term P6 mission
  concepts.
\end{enumerate}

\end{document}
